\documentclass[11pt,a4paper]{article}
\usepackage[a4paper,margin=2cm]{geometry}
\usepackage[T1]{fontenc}
\usepackage[utf8]{inputenc}
\usepackage{lmodern}
\usepackage{xcolor}
\usepackage{tabularx}
\usepackage{booktabs}
\usepackage{enumitem}
\usepackage{titlesec}
\usepackage{fancyhdr}
\usepackage[hidelinks]{hyperref}
\usepackage{tcolorbox}
\usepackage{array}
\usepackage{amssymb}

\definecolor{phgreen}{HTML}{0B6E4F}
\definecolor{phlight}{HTML}{E8F5EF}
\definecolor{phgrey}{HTML}{55606B}

\renewcommand{\familydefault}{\sfdefault}
\setlength{\parindent}{0pt}
\setlength{\parskip}{6pt}
\setlist{nosep,leftmargin=1.4em}

\titleformat{\section}{\Large\bfseries\color{phgreen}}{}{0pt}{}
\titlespacing*{\section}{0pt}{14pt}{6pt}
\titleformat{\subsection}{\large\bfseries\color{phgrey}}{}{0pt}{}
\titlespacing*{\subsection}{0pt}{10pt}{4pt}

\pagestyle{fancy}
\fancyhf{}
\fancyhead[L]{\small\color{phgrey}Public Health Colosseum --- Admin Guide}
\fancyhead[R]{\small\color{phgrey}Version 1}
\fancyfoot[C]{\small\color{phgrey}Powered by SCIAT SOLUTIONS \quad·\quad \thepage}
\renewcommand{\headrulewidth}{0.4pt}

\newcommand{\cmd}[1]{\texttt{\textbf{#1}}}
\newtcolorbox{tip}{colback=phlight,colframe=phgreen,boxrule=0.6pt,arc=2pt,left=6pt,right=6pt,top=4pt,bottom=4pt}

\newenvironment{commands}{%
  \renewcommand{\arraystretch}{1.35}%
  \tabularx{\textwidth}{@{}>{\raggedright\arraybackslash}p{5.6cm}X@{}}\toprule
  \textbf{Command} & \textbf{What it does} \\\midrule}
  {\bottomrule\endtabularx}
\newenvironment{settings}{%
  \renewcommand{\arraystretch}{1.35}%
  \tabularx{\textwidth}{@{}>{\raggedright\arraybackslash}p{5.0cm}>{\raggedright\arraybackslash}X>{\raggedright\arraybackslash}p{3.8cm}@{}}\toprule
  \textbf{Setting} & \textbf{What it is for} & \textbf{Where to get it} \\\midrule}
  {\bottomrule\endtabularx}
\newcommand{\shell}[1]{\begin{tcolorbox}[colback=black!85,colframe=black!85,coltext=white,boxrule=0pt,arc=2pt,left=6pt,top=3pt,bottom=3pt]\ttfamily #1\end{tcolorbox}}


\begin{document}

\begin{center}
{\Huge\bfseries\color{phgreen} Public Health Colosseum}\\[4pt]
{\Large Admin Guide --- Running and Managing the Bot}\\[2pt]
{\color{phgrey} Server commands, API keys and settings, WhatsApp admin commands}
\end{center}

\section{Before you start}

Log in to the server over SSH and go to the bot's folder. Every server command below is run from there.

\shell{cd /root/phc-bot}

\textbf{First install on a new server} (one command, then answer the questions):

\shell{curl -fsSL https://raw.githubusercontent.com/Maerih/phc-bot/main/install.sh | bash}

\section{Running the bot}

\begin{commands}
\cmd{./ctl.sh start} & Start everything: checks the server, Docker, settings, builds, database, bot, and shows the QR code if WhatsApp needs pairing. Safe to run any time --- it only changes what needs changing. \\
\cmd{./ctl.sh stop} & Stop the bot and the database \\
\cmd{./ctl.sh restart} & Restart the bot (keeps the WhatsApp session) \\
\cmd{./ctl.sh status} & Is it running, and is WhatsApp connected? \\
\cmd{./ctl.sh doctor} & Run every check, including that each API key works. Changes nothing. \textbf{Run this first when anything is wrong.} \\
\cmd{./ctl.sh logs} & Watch the live log (Ctrl-C stops watching) \\
\cmd{./ctl.sh update} & Get the latest version from GitHub and restart. Keeps data, settings and session. \\
\end{commands}

\section{WhatsApp connection}

\begin{commands}
\cmd{./ctl.sh pair} & Show the QR code in the terminal. Scan with WhatsApp $\rightarrow$ Settings $\rightarrow$ Linked devices $\rightarrow$ Link a device. \\
\cmd{./ctl.sh reset-session} & Unlink and pair again with a fresh QR code. Use when the bot is connected but does not answer. Data and settings are kept. \\
\end{commands}

\begin{tip}
\textbf{The one rule:} never run the bot on two machines with the same WhatsApp session --- it breaks the session. Moving servers? Pair the new one fresh, then remove the old device in WhatsApp's Linked devices.
\end{tip}

\section{API keys and settings}

\subsection{The easy way: answer the questions}

\shell{./ctl.sh setup \\ ./ctl.sh start}

Setup asks for the admin numbers and every key. Keys are hidden as you type; pressing Enter keeps the current value. Then \cmd{./ctl.sh start} applies the changes (the bot restarts only if something changed).

\subsection{Or edit the settings file directly}

\shell{nano .env \\ ./ctl.sh start}

In nano: change the line, then Ctrl-O, Enter to save, Ctrl-X to exit. Each setting is one line: \cmd{NAME=value}, no spaces around the \cmd{=}.

\subsection{The settings}

\begin{settings}
\cmd{ADMIN\_NUMBERS} & Admin WhatsApp numbers, country code first, commas between: \cmd{254712345678,254798765432} & --- \\
\cmd{DEEPSEEK\_API\_KEY} & The AI: \cmd{.ask}, AI search, choosing which jobs go out in alerts, and the CV and application assistant (needed for it) & platform.deepseek.com/ api\_keys \\
\cmd{TAVILY\_API\_KEY} & Web search for \cmd{.ask} and the search commands. Also set \cmd{WEB\_SEARCH\_PROVIDER=tavily} (setup does this). & app.tavily.com (free plan) \\
\cmd{JSEARCH\_API\_KEY} & LinkedIn and other big job sites, searched daily & rapidapi.com $\rightarrow$ JSearch $\rightarrow$ free plan \\
\cmd{RELIEFWEB\_APPNAME} & ReliefWeb NGO, UN and humanitarian jobs & apidoc.reliefweb.int \\
\cmd{AI\_PROVIDER} & \cmd{auto} (DeepSeek first), \cmd{deepseek}, \cmd{gemini}, \cmd{anthropic} or \cmd{none} & --- \\
\cmd{GEMINI\_API\_KEY}, \cmd{ANTHROPIC\_API\_KEY} & Optional back-up AIs, used only if DeepSeek has no key & aistudio.google.com, console.anthropic.com \\
\cmd{TRACKER\_ENABLED} & \cmd{on} or \cmd{off}: the opportunity alerts & --- \\
\cmd{TRACKER\_INTERVAL\_MINUTES} & How often the job sites are checked (default 30) & --- \\
\cmd{TIMEZONE} & For the alert schedule and deadlines (default \cmd{Africa/Nairobi}) & --- \\
\cmd{COMMAND\_PREFIX} & The character commands start with (default \cmd{.}) & --- \\
\end{settings}

After any change, \cmd{./ctl.sh doctor} shows each key as working (\textcolor{phgreen}{\checkmark}) or what is wrong with it.

\section{Data and safety}

\begin{commands}
\cmd{./ctl.sh backup} & Save the database and WhatsApp session into \cmd{backups/}. Keep these files private: they can link a device to your WhatsApp. \\
\cmd{./ctl.sh autostart on} & Start the bot by itself after the server reboots (\cmd{off} to undo) \\
\cmd{./ctl.sh use-system-docker} & Move the bot, with its database, onto the server's own Docker. Only when \cmd{doctor} says the bot cannot reach the internet. \\
\cmd{./ctl.sh uninstall} & Remove the bot, its image and its \textbf{database}. The WhatsApp session is kept. Take a backup first. \\
\end{commands}

\section{Admin commands in WhatsApp}

Sent from a number in \cmd{ADMIN\_NUMBERS}:

\begin{commands}
\cmd{.alerts on} / \cmd{.alerts off} & Opportunity alerts in this chat (group admins can too) \\
\cmd{.alerts} & Status: sites watched, last check, how many opportunities are waiting \\
\cmd{.scan} & Check every opportunity site right now and report what was found \\
\cmd{.pushall} & Send everything waiting into this chat at once, as a list and a spreadsheet \\
\cmd{.alerts interval 3-5} & Gap between alert cards: here a random 3 to 5 minutes; any 1--120 \\
\cmd{.alerts pause} / \cmd{.alerts resume} & Stop or restart alerts in every chat at once \\
\cmd{.ping} & Is the bot alive? \\
\end{commands}

\subsection{Bot number switches}

All are \textbf{off} until you turn them on, with \cmd{on} or \cmd{off} after the name. They act for the bot's own WhatsApp number.

\begin{commands}
\cmd{.autoview on} & View contacts' status updates automatically \\
\cmd{.autolike on} & React \textcolor{phgreen}{\(\heartsuit\)} to contacts' status updates \\
\cmd{.alwaysonline on} & Show the bot as online all the time \\
\cmd{.recording on} & Show ``recording audio\ldots'' instead of ``typing\ldots'' while the bot works on a reply \\
\cmd{.contacts} & Everyone who has messaged the bot, as one contacts file to save on a phone in one tap \\
\end{commands}

\begin{tip}
\textbf{Use sparingly.} WhatsApp can limit numbers that look automated; \cmd{alwaysonline} and \cmd{autolike} are the most noticeable. Saving deleted messages is deliberately not offered: members delete messages on purpose.
\end{tip}

\subsection{AI in private chats (personal number)}

If the bot runs on your own WhatsApp number, decide when the AI may answer private messages. Commands such as \cmd{.cv} always work for everyone; this is only about ordinary messages.

\begin{commands}
\cmd{.aimode off} & \textbf{The default.} Your private messages are yours. The AI answers only someone who starts with an AI command (\cmd{.cv}, \cmd{.ask}, \cmd{.profile}, \cmd{.apply}\ldots), and keeps talking with them for 30 minutes after their last message \\
\cmd{.aimode on} & The AI answers every private message \\
\cmd{.aimode} & Which mode is on, and how many conversations are open \\
\cmd{.aimode stop} \textit{(in a chat)} & The AI never answers in that chat; \cmd{.aimode start} to undo. In someone else's chat the confirmation goes to your own ``Message yourself'' chat, not to them \\
\cmd{.stop} & End the AI conversation in this chat right now. Members can use it too \\
\end{commands}

\textbf{Taking over:} as soon as you type in a chat yourself, the AI stops answering there (for 12 hours when the mode is on). Your ``Message yourself'' chat follows the same rules: start with \cmd{.ask} or \cmd{.cv} to talk to the AI there, and \cmd{.stop} when you are done.

\subsection{Career assistant}

Members chat privately with the bot to build CVs and prepare applications; see the User Guide. It needs \cmd{DEEPSEEK\_API\_KEY}; \cmd{TAVILY\_API\_KEY} lets it find the official application page. CVs are made as PDF and Word files inside the bot's container (it includes a browser for this), stored in the database, and deleted together with the member's profile after 12 months without use or when they send \cmd{delete my data}. Email applications are prepared for the member to send; simple online forms are filled only after the member checks a screenshot and replies \cmd{confirm}. No mailbox or passwords are stored.

\textbf{Alert schedule (Nairobi time):} Monday 3:00 AM -- Friday 1:00 PM, a card every few minutes; Friday 2:00 PM the weekly report and spreadsheet, then the weekend message; nothing at the weekend. Group-management and moderation commands are in the User Guide.

\section{When something is wrong}

\begin{enumerate}[label=\textbf{\arabic*.},leftmargin=1.6em,itemsep=3pt]
  \item Run \cmd{./ctl.sh doctor} and read the lines marked \textbf{$\times$} --- each says how to fix it.
  \item \textbf{Connected but not answering:} \cmd{./ctl.sh reset-session}, then scan the QR again.
  \item \textbf{After installing other software on the server:} \cmd{./ctl.sh update}, then \cmd{./ctl.sh doctor}.
  \item \textbf{AI or web answers missing:} \cmd{./ctl.sh doctor} shows which key is missing or rejected; add it with \cmd{./ctl.sh setup}.
  \item \textbf{No alerts arriving:} send \cmd{.alerts} --- check it says ON and something is waiting; remember the weekend pause.
  \item \textbf{Still stuck:} share the output of \cmd{./ctl.sh doctor} and a minute of \cmd{./ctl.sh logs}.
\end{enumerate}

\vfill
\begin{center}
\color{phgrey}\small
\textbf{Public Health Colosseum} · Powered by \textbf{SCIAT SOLUTIONS}
\end{center}

\end{document}
